

\chapter{Khi Thầy Cô Cũng Lý Sự}
\markboth{Khi Thầy Cô Cũng Lý Sự}{Khi Thầy Cô Cũng Lý Sự}

\section{``Thầy Dạy Vậy Là Đúng Rồi, Tra Sách Làm Gì''}
\begin{truyen}{``Thầy Dạy Vậy Là Đúng Rồi, Tra Sách Làm Gì''}{``My Way Is Correct. Why Check the Book?''}
\chuhoa{H}{ọc sinh:} ``Thầy ơi, em tra sách thấy cách giải khác với cách thầy dạy. Cách nào đúng ạ?''
\textit{(Student: ``Teacher, the textbook shows a different method from what you taught. Which one is correct?'')}

\teacherpair{Cách thầy dạy là cách thầy thấy hợp lý nhất. Em không cần phải tra thêm.}{``The method I teach is what I consider most appropriate. You do not need to look further.''}

Học sinh im, nhưng trong đầu vẫn bấm máy tính thử lại. Cả hai cách đều đúng. Cách sách hợp lý hơn.
\textit{(The student goes quiet, but mentally rechecks with a calculator. Both methods work. The textbook approach is cleaner.)}

\begin{baihoc}
Một giáo viên giỏi không sợ khi học sinh tìm thấy cách khác. Họ sẽ hỏi ngược lại: ``Em thấy điểm khác ở đâu? Cách đó tốt không?''
\textit{(A good teacher is not threatened when a student finds a different method. They ask instead: ``Where do you see the difference? Is that approach better?'')}

Sợ bị bắt bẻ là dấu hiệu của sự không tự tin, không phải của uy quyền.
\textit{(Fear of being challenged is a sign of insecurity, not authority.)}
\end{baihoc}
\end{truyen}

\section{``Em Hỏi Nhiều Quá, Phá Lớp Học''}
\begin{truyen}{``Em Hỏi Nhiều Quá, Phá Lớp Học''}{``You Ask Too Much. You Are Disrupting the Class''}
\chuhoa{H}{ọc sinh:} ``Thầy ơi tại sao phải làm bước này? Bỏ qua được không ạ?''
\textit{(Student: ``Why do we need this step, teacher? Can it be skipped?'')}

\teacherpair{Thôi đừng hỏi linh tinh nữa. Cứ làm theo là xong.}{``Stop asking unnecessary things. Just follow along.''}

Cả lớp im. Đứa học sinh ghi chép nhưng không hiểu tại sao. Ba tuần sau kiểm tra, nó làm sai y chang bước đó vì chưa bao giờ được giải thích.
\textit{(The class falls silent. The student copies notes but does not understand why. Three weeks later at the test, they make the same mistake in that exact step because it was never explained.)}

\begin{baihoc}
Câu hỏi ``tại sao'' là dấu hiệu của tư duy. Dập tắt nó không phải dạy học, đó là kiểm soát.
\textit{(The question ``why'' is a sign of thinking. Silencing it is not teaching. It is control.)}

Học sinh hỏi nhiều thường không phải học sinh hỗn. Đó là học sinh đang cố hiểu thật sự.
\textit{(Students who ask a lot are usually not troublemakers. They are students who are genuinely trying to understand.)}
\end{baihoc}
\end{truyen}

\section{Đúng Về Nội Dung, Sai Về Cách}
\begin{truyen}{Đúng Về Nội Dung, Sai Về Cách}{Right in Content, Wrong in Delivery}
\chuhoa{H}{ọc sinh:} ``Thầy hay mắng học sinh trước mặt cả lớp. Tụi em xấu hổ lắm.''
\textit{(Student: ``You scold students in front of the whole class. It is very humiliating for us.'')}

\teacherpair{Thầy mắng vì muốn em tiến bộ. Em không biết ơn thì thôi, chứ đừng có phàn nàn.}{``I scold because I want you to improve. If you are not grateful, fine. But do not complain.''}

Một trò khác lén ghi vào sổ tay: \textit{``Thầy đúng về nội dung. Sai về cách. Và không nhận ra sự khác biệt đó.''}
\textit{(Another student secretly writes in a notebook: ``The teacher is right about the content. Wrong about the method. And does not recognize the difference.'')}

\begin{baihoc}
Ý tốt không miễn trừ khỏi trách nhiệm với phương pháp. Nhục mạ công khai không bao giờ là công cụ giáo dục, dù nội dung phê bình hoàn toàn chính xác.
\textit{(Good intent does not exempt responsibility for method. Public humiliation is never a teaching tool, even when the criticism is entirely accurate.)}
\end{baihoc}
\end{truyen}

\section{``Đừng Cãi Thầy, Thầy Kinh Nghiệm Hơn Em''}
\begin{truyen}{``Đừng Cãi Thầy, Thầy Kinh Nghiệm Hơn Em''}{``Do Not Argue With Me. I Have More Experience''}
\chuhoa{H}{ọc sinh:} ``Nhưng thưa thầy, trong tình huống đó, cách em làm cũng hợp lý mà ạ...''
\textit{(Student: ``But teacher, in that situation, my approach was also reasonable...'')}

\teacherpair{Thầy dạy hai mươi năm rồi. Em học mấy năm? Biết mình là ai chưa?}{``I have been teaching for twenty years. How many years have you studied? Know your place.''}

Học sinh cúi đầu. Im lặng. Nhưng trong nhóm chat tối hôm đó: \textit{``Mình hỏi mấy anh khoá trên rồi, cách mình đúng hơn thật.''} 
\textit{(The student bows their head. Silence. But in the class chat that evening: ``I asked seniors and they confirmed: my method was actually more correct.'')}

\begin{baihoc}
Kinh nghiệm là tài sản quý. Nhưng kinh nghiệm không phải là chiếc khiên để né mọi câu hỏi. Giáo viên cứng đầu mà sai thì nguy hiểm gấp đôi, vì họ có uy quyền để khiến sai trở thành đúng trong mắt học sinh.
\textit{(Experience is a valuable asset. But experience is not a shield against every question. A stubborn teacher who is wrong is doubly dangerous, because they have the authority to make the class accept wrong as right.)}
\end{baihoc}
\end{truyen}

\section{Chấm Điểm Theo Cảm Tình}
\begin{truyen}{Chấm Điểm Theo Cảm Tình}{Grading by Personal Feeling}
\chuhoa{H}{ọc sinh:} ``Tụi em thấy bài em A điểm thấp hơn bài em B dù viết tương đương, chỉ khác là em B hay phát biểu.''
\textit{(Student: ``We noticed that student A got a lower score than student B despite similar writing. The only difference is that student B participates in class often.'')}

\teacherpair{Em không hiểu cách chấm của thầy. Chấm bài là phức tạp lắm.}{``You do not understand how I grade. Grading is complicated.''}

Không có lời giải thích thêm. Chỉ có câu trả lời đóng cửa mọi thắc mắc thêm.
\textit{(No further explanation. Just an answer that closes every follow-up question.)}

\begin{baihoc}
Tiêu chí chấm điểm không rõ ràng là mảnh đất màu mỡ cho sự thiên vị. Học sinh không yêu cầu hoàn hảo. Họ chỉ yêu cầu được biết mình bị đánh giá theo cái gì.
\textit{(Unclear grading criteria are fertile ground for bias. Students do not ask for perfection. They only ask to know what they are being judged on.)}
\end{baihoc}
\end{truyen}

\section{``Em Nghĩ Gì Cũng Được, Miễn Đừng Nói Trước Mặt Thầy''}
\begin{truyen}{``Em Nghĩ Gì Cũng Được, Miễn Đừng Nói Trước Mặt Thầy''}{``Think Whatever You Want, Just Not in Front of Me''}
\chuhoa{H}{ọc sinh:} ``Thầy bao giờ sai chưa ạ?''
\textit{(Student: ``Have you ever been wrong, teacher?'')}

\teacherpair{Sinh ra đã không sai. Chứ thầy cũng sai. Nhưng không phải chuyện của em.}{``Born perfect? No. I have been wrong. But that is not your business.''}

Nhiều học sinh trao nhau ánh mắt. Câu trả lời đó thực ra thừa nhận sai, nhưng lại từ chối cơ hội cho học sinh thấy rằng thừa nhận sai là điều bình thường và dũng cảm.
\textit{(Several students exchange glances. The answer technically admits fault, but refuses the chance to show students that admitting mistakes is normal and courageous.)}

\begin{baihoc}
Giáo viên thừa nhận sai trước học sinh không mất uy. Ngược lại, họ trở thành hình mẫu của điều quan trọng nhất mà trường học cần dạy: sự trung thực với bản thân.
\textit{(A teacher who admits a mistake in front of students does not lose authority. On the contrary, they become a model of the most important thing a school should teach: honesty with oneself.)}
\end{baihoc}
\end{truyen}

\section{Giao Bài Tập Không Hỏi Học Sinh Đang Gánh Gì}
\begin{truyen}{Giao Bài Tập Không Hỏi Học Sinh Đang Gánh Gì}{Assigning Homework Without Asking What Students Already Carry}
\chuhoa{H}{ọc sinh:} ``Thầy ơi tuần này tụi em thi đủ thứ mà thầy giao thêm hai bài tập lớn.''
\textit{(Student: ``This week we have exams in almost every subject, but you still assigned two large projects.'')}

\teacherpair{Phần thầy không liên quan đến phần môn khác. Môn thầy thì thầy quyết.}{``My subject is not connected to other subjects. In my class, I decide.''}

Đúng, thầy có quyền quyết. Nhưng học sinh thì không có quyền phân thân.
\textit{(Technically, the teacher has the right to decide. But students do not have the ability to split themselves.)}

\begin{baihoc}
Hệ thống giáo dục thiếu điều phối ngang giữa giáo viên các môn. Hậu quả là học sinh đôi khi chết đuối không phải vì một môn mà vì tất cả đồng loạt dồn vào cùng một tuần.

Người dạy có thể không có lỗi hệ thống, nhưng họ có lựa chọn: cứng nhắc hay linh hoạt.
\textit{(The education system lacks horizontal coordination among teachers. Students sometimes drown not from one subject but from all subjects piling on the same week. A teacher may not be responsible for the system's flaw, but they do have a choice: rigid or flexible.)}
\end{baihoc}
\end{truyen}

\begin{baihoc}
Khi giáo viên cũng cùn, thiệt hại nặng hơn học sinh cùn --- vì giáo viên có quyền lực, có ảnh hưởng lâu dài, và thường ít bị ai bắt bẻ hơn.
\textit{(When teachers are stubborn, the damage is greater than when students are stubborn, because teachers hold power, create lasting influence, and are rarely challenged.)}

Phòng học không phải nơi chỉ có một phía cần tự soi.
\textit{(The classroom is not a place where only one side needs self-reflection.)}
\end{baihoc}

\begin{ghinhoanh}
\textbf{Short English to remember:}

Power without humility is a classroom hazard.

The best teachers are still learning.
\end{ghinhoanh}

\begin{tuhocnhanh}
1. Em đã từng gặp giáo viên cùn chưa? Phản ứng của em khi đó là gì? \textit{(Have you ever encountered a stubborn teacher? What was your reaction?)}

2. Nếu là giáo viên, em sẽ làm gì khi học sinh bắt bẻ đúng điểm mình sai? \textit{(If you were a teacher, what would you do when a student correctly catches your mistake?)}
\end{tuhocnhanh}

\ngancach
